% Appendix A v4 -- BGFM (Boltzmann-Guided Flow Matching) formulation.
%
% Supersedes appendix_A_v3.tex (bond-free physics-drift formulation).
% Adds the three-term training objective and proves Boltzmann consistency.
%
% Theorems 1-5 carry over from v3 (manifold geometry + tangent projection
% + retracted Euler convergence). New Theorems 6-8 cover the BGFM
% objective: score formula, force-consistency equivalence, and Boltzmann
% consistency in the large-lambda limit.

\documentclass[11pt]{article}
\usepackage{amsmath,amssymb,amsthm}
\usepackage[margin=1in]{geometry}
\usepackage{mathtools}

\newtheorem{theorem}{Theorem}
\newtheorem{lemma}[theorem]{Lemma}
\newtheorem{proposition}[theorem]{Proposition}
\newtheorem{corollary}[theorem]{Corollary}
\newtheorem{definition}[theorem]{Definition}
\theoremstyle{remark}
\newtheorem{remark}[theorem]{Remark}
\newtheorem{assumption}[theorem]{Assumption}
\theoremstyle{plain}

\newcommand{\M}{\mathcal{M}}
\newcommand{\X}{\mathcal{X}}
\newcommand{\A}{\mathcal{A}}
\newcommand{\C}{\mathcal{C}}
\newcommand{\R}{\mathbb{R}}
\newcommand{\E}{\mathbb{E}}
\newcommand{\T}{\mathsf{T}}
\newcommand{\ster}{\mathrm{ster}}
\newcommand{\FM}{\mathrm{FM}}
\newcommand{\force}{\mathrm{force}}
\newcommand{\energy}{\mathrm{energy}}
\newcommand{\TV}{\mathrm{TV}}
\newcommand{\KL}{\mathrm{KL}}
\newcommand{\BGFM}{\mathrm{BGFM}}
\newcommand{\Var}{\mathrm{Var}}
\newcommand{\norm}[1]{\left\|#1\right\|}
\newcommand{\ip}[2]{\langle #1, #2\rangle}

\title{Appendix A v4: BGFM Proofs \\ \large Boltzmann-Guided Flow Matching: Score, Force, Energy Alignment}
\author{(iterating -- matches notes/bgfm_method.md)}
\date{}

\begin{document}
\maketitle

This appendix extends the v3 formulation (bond-free physics-drift flow
matching) with the three-term BGFM training objective. We prove that the
learned density converges to the Boltzmann distribution under a neural
potential $E$ in the limit of large energy-consistency weight.

Sections 1--3 carry over verbatim from v3 (setup, steric manifold,
tangent projection/retraction). We focus here on the new BGFM-specific
content: Sections 4 (score formula), 5 (force consistency), 6 (energy
consistency), and the main theorem in Section 7.

\section{Recap: setup and notation}

From v3 (condensed). Configurations live in
$\X = \R^{3N} \times (\Delta^{A-1})^N \times (\Delta^{C-1})^N$, with
steric-admissible manifold $\M_\ster = \{x : \norm{r_i - r_j}^2 \ge
d_\mathrm{eff}^2(\tilde a_i, \tilde a_j) \; \forall i \ne j\}$.

A universal neural potential $E : \X \to \R$ supplies the Boltzmann
reference
\[
p_\mathrm{Boltz}(x) := Z^{-1} \exp\!\big(-E(x)/kT\big).
\]
At each training point $(x, F)$ where $x \sim p_\text{data}$ and $F(x)
= -\nabla_r E(x)$ is the DFT force (precomputed offline), the BGFM
objective enforces three alignments of the model-induced density
$p_\theta$:

\begin{itemize}
\item \textbf{FM (data).} $p_\theta$ assigns density to observed data.
\item \textbf{Force (local gradient).} $\nabla \log p_\theta(x)
      \approx F(x)/kT$ at data points.
\item \textbf{Energy (global log-density shape).} $\log p_\theta(x) +
      E(x)/kT$ is constant across data points.
\end{itemize}

\section{Score of a flow-matching marginal density (recap + details)}

\begin{lemma}[FM score formula, Gaussian prior]\label{lem:score}
Let $p_0 = \mathcal{N}(0, \sigma^2 I)$ and let $u^*(x, t) = \E[x_1 - x_0
\mid x_t = x]$ be the marginal velocity of the linear interpolant
$x_t = (1 - t) x_0 + t x_1$. For $t \in (0, 1)$, the marginal score is
\begin{equation}\label{eq:score}
s^*(x, t) \;:=\; \nabla_x \log p_t(x) \;=\; \frac{t \cdot u^*(x, t) - x}{(1 - t)\, \sigma^2}.
\end{equation}
\end{lemma}

\begin{proof}
The conditional $p_t(x \mid x_1) = \mathcal{N}(t x_1, (1-t)^2 \sigma^2 I)$
so $\nabla_x \log p_t(x \mid x_1) = -(x - t x_1) / ((1-t)^2 \sigma^2)$.
Marginalizing over $x_1$ with the conditional law $p(x_1 \mid x_t = x)$:
\begin{align*}
s^*(x, t)
  &= \E[\nabla_x \log p_t(x \mid x_1) \mid x_t = x] \\
  &= -(x - t\, \E[x_1 \mid x_t = x]) / ((1-t)^2 \sigma^2).
\end{align*}
From $x_t = (1 - t) x_0 + t x_1$, solving for $x_1$:
$x_1 = (x_t - (1-t) x_0)/t$. Thus
$\E[x_1 \mid x_t = x] = x + (1-t) \E[-x_0 \mid x_t = x] / t + \cdots$,
and using $u^*(x, t) = \E[x_1 - x_0 \mid x_t = x]$ together with
$x = (1-t) \E[x_0 \mid x_t = x] + t \E[x_1 \mid x_t = x]$, we obtain
$\E[x_1 \mid x_t = x] = x + (1-t) u^*(x, t)$. Substituting:
\begin{align*}
s^*(x, t) &= -(x - t(x + (1-t) u^*)) / ((1-t)^2 \sigma^2) \\
          &= -((1-t) x - t (1-t) u^*) / ((1-t)^2 \sigma^2) \\
          &= (t u^* - x) / ((1-t) \sigma^2).
\end{align*}
\end{proof}

\begin{remark}[Implementation]
Equation~\eqref{eq:score} is used in \texttt{cfm\_mol.bgfm\_loss.
score\_from\_fm\_velocity}, evaluated at $t_\text{eval}$ slightly below
$1$ (default $0.95$) to avoid the $1/(1-t)$ singularity. Unit tests
against a 1D Gaussian target confirm the formula to numerical precision
(\texttt{tests/test\_bgfm\_loss.py::test\_score\_from\_fm\_velocity\_one\_d\_gaussian}).
\end{remark}

\section{Force consistency}

\begin{definition}[Force-consistency loss]
Given training data $\{(x^{(m)}, F^{(m)})\}_{m=1}^M$ with $F^{(m)}
= -\nabla_r E(x^{(m)})$, define
\begin{equation}\label{eq:lforce}
\mathcal{L}_\force(\theta)
  = \frac{1}{M} \sum_{m=1}^M
       \norm{s_\theta(x^{(m)}, t_\mathrm{eval}) - F^{(m)} / kT}^2,
\end{equation}
where $s_\theta(\cdot, t) = (t \cdot v_\theta - x)/((1-t)\sigma^2)$ is
the FM-implied score (Lemma~\ref{lem:score}).
\end{definition}

\begin{proposition}[Force-consistency $\Leftrightarrow$ local Boltzmann
score]\label{prop:force-equiv}
If $\mathcal{L}_\force(\theta) = 0$ (achievable at infinite capacity),
then for every $x \in \mathrm{supp}(p_\text{data})$,
\[
\nabla_x \log p_\theta(x, t_\mathrm{eval}) = -\nabla_x E(x) / kT.
\]
In the $t_\mathrm{eval} \to 1$ limit (so $p_\theta(\cdot, 1) = p_\theta$
is the terminal model distribution), we have
$\nabla \log p_\theta(x) = -\nabla E(x)/kT$ on data support.
\end{proposition}

\begin{proof}
Direct substitution of $\mathcal{L}_\force = 0$ into~\eqref{eq:lforce}
yields $s_\theta(x, t_\mathrm{eval}) = F(x)/kT = -\nabla E(x)/kT$ at
each training point, pointwise on support. Continuity of $s_\theta$
in $t$ extends this to $t \to 1$.
\end{proof}

\begin{remark}[``Local Boltzmann''
consistency]\label{rmk:local-boltz}
Proposition~\ref{prop:force-equiv} gives that $\log p_\theta$ has
the correct \emph{gradient} matching Boltzmann at every data point.
This determines $\log p_\theta$ on the data support \emph{up to an
additive constant}. The energy term (Section 5) fixes this constant.
\end{remark}

\section{Energy consistency via variance form}

\begin{definition}[Energy-consistency loss]
Given training data with precomputed DFT energies $\{E^{(m)}\}_{m=1}^M$
and the model's log-density $\log p_\theta$ (computed via
change-of-variables, Eq.~5.1 of main text):
\begin{equation}\label{eq:lenergy}
\mathcal{L}_\energy(\theta)
  = \Var_{m \sim \mathrm{Unif}[M]}\!\big[\log p_\theta(x^{(m)}) + E^{(m)}/kT\big].
\end{equation}
\end{definition}

\begin{remark}[Why variance, not squared error?]
The target of the $\log p$ equation is
$\log p_\theta(x) = -E(x)/kT - \log Z$, where $Z = \int e^{-E/kT} dx$
is unknown and notoriously hard to estimate. A squared-error loss
$(\log p_\theta + E/kT + \log Z)^2$ would require learning or
estimating $\log Z$. The variance formulation~\eqref{eq:lenergy} is
\emph{invariant under constant shifts}: for any $c \in \R$, $\Var(X + c)
= \Var(X)$. Thus $\log Z$ is automatically absorbed and we need only
enforce the \emph{shape} of $\log p_\theta$ to match $-E/kT$.
\end{remark}

\begin{proposition}[Energy-consistency $\Leftrightarrow$ global
Boltzmann shape]\label{prop:energy-equiv}
$\mathcal{L}_\energy(\theta) = 0$ iff there exists $c \in \R$ such that
$\log p_\theta(x) + E(x)/kT = c$ for all $x \in \mathrm{supp}(p_\text{data})$.
Equivalently, $p_\theta(x) = e^{-c} \exp(-E(x)/kT)$ on support.
\end{proposition}

\begin{proof}
$\Var(X) = 0$ iff $X$ is almost-surely constant. Apply to $X = \log
p_\theta + E/kT$ evaluated on training samples.
\end{proof}

\section{Full BGFM objective and Boltzmann consistency}

\begin{definition}[BGFM training objective]
\begin{equation}\label{eq:bgfm-total}
\mathcal{L}_\BGFM(\theta)
  = \mathcal{L}_\FM(\theta)
  + \lambda_1\, \mathcal{L}_\force(\theta)
  + \lambda_2\, \mathcal{L}_\energy(\theta),
\end{equation}
with weights $\lambda_1, \lambda_2 \ge 0$ scheduled to ramp from $0$
(warmup) to $\lambda_1^\star, \lambda_2^\star$ (full value) during
training.
\end{definition}

\begin{assumption}[BGFM regularity]\label{assum:bgfm}
\begin{enumerate}
\item[(B1)] The FM loss achieves $\mathcal{L}_\FM(\theta^*) \le \varepsilon_\FM$
      at the optimum, ensuring $p_\theta(\cdot, 1)$ covers
      $\mathrm{supp}(p_\text{data})$.
\item[(B2)] The model family has sufficient capacity for
      $v_\theta$ to interpolate $u^*$ and express any Boltzmann density
      under~$E$.
\item[(B3)] The neural potential $E$ is $C^1$ with $L_E$-Lipschitz
      gradient and $E_{\min} := \inf E > -\infty$.
\end{enumerate}
\end{assumption}

\begin{theorem}[Boltzmann consistency]\label{thm:boltzmann}
Under Assumption~\ref{assum:bgfm}, let $\theta^\star(\lambda_2)$
minimize~\eqref{eq:bgfm-total} with $\lambda_1 > 0$ fixed and
$\lambda_2 > 0$ growing. Then as $\lambda_2 \to \infty$, the induced
density satisfies, for every $x \in \mathrm{supp}(p_\text{data})$:
\[
p_{\theta^\star(\lambda_2)}(x) \;\longrightarrow\; \frac{\exp(-E(x)/kT)}{Z^\star},
\qquad Z^\star = \int_{\mathrm{supp}(p_\text{data})} \exp(-E(x')/kT)\, dx'.
\]
\end{theorem}

\begin{proof}
\emph{Step 1 (energy loss forces Boltzmann shape).} As $\lambda_2 \to
\infty$, minimizing~\eqref{eq:bgfm-total} forces $\mathcal{L}_\energy
\to 0$. By Proposition~\ref{prop:energy-equiv}, this is equivalent to
$\log p_{\theta^\star} + E/kT$ being constant on $\mathrm{supp}(p_\text{data})$.

\emph{Step 2 (FM loss pins the support).} By Assumption~(B1),
$p_{\theta^\star}$ covers $\mathrm{supp}(p_\text{data})$, with error
$\varepsilon_\FM$ in a suitable metric (e.g., reverse KL).

\emph{Step 3 (normalization).} The constant $c$ in
Proposition~\ref{prop:energy-equiv} is determined by $\int p_{\theta^\star} = 1$
on support. Since $p_{\theta^\star}(x) = e^{-c} \exp(-E(x)/kT)$ on
support, integration gives $e^{-c} = 1/Z^\star$ where $Z^\star$ is as
stated.

\emph{Step 4 (force loss as consistency check).} Proposition~\ref{prop:force-equiv}
with $\lambda_1 > 0$ fixed ensures that the local gradient of
$\log p_{\theta^\star}$ matches $-\nabla E/kT$, which is automatically
true under Step 1 (by differentiation). Hence $\lambda_1$ is redundant
in the limit but accelerates convergence at finite $\lambda_2$ by
providing a parallel, local supervision signal.

Combining: on $\mathrm{supp}(p_\text{data})$, $p_{\theta^\star}(x) =
\exp(-E(x)/kT)/Z^\star$.
\end{proof}

\begin{corollary}[Chemistry-agnostic Boltzmann sampling]\label{cor:chem-free-boltz}
Theorem~\ref{thm:boltzmann} holds for any chemistry covered by $E$'s
training domain. With $E = E_\text{OMol25}$ (83-element, ωB97M-V
accuracy), the model's terminal distribution $p_{\theta^\star}$ is
Boltzmann-aligned across organics, organometallics, radicals, and
hypervalent species uniformly -- no chemistry-specific valence table,
bond-order constraint, or rule-based sanitization is required.
\end{corollary}

\begin{remark}[Why three terms rather than one?]
In the infinite-capacity, infinite-data limit, the three BGFM terms are
redundant by the chain of equivalences (FM $\Rightarrow$ support,
energy $\Rightarrow$ shape, force = gradient of energy). In practice
each provides a different training signal with different learning
dynamics:

\begin{itemize}
\item $\mathcal{L}_\FM$ is a \emph{velocity-field} loss (dense, per-sample,
      fast). It establishes the model's coverage before energy/force can
      be meaningfully applied.
\item $\mathcal{L}_\force$ is a \emph{gradient} loss (per-atom, per-sample,
      medium speed). It sharpens the local shape without requiring
      full trajectory integration.
\item $\mathcal{L}_\energy$ is a \emph{global} loss requiring $\log p_\theta$
      via trajectory divergence integral (per-graph, per-sample, slow).
      It constrains the overall density shape but is expensive.
\end{itemize}

Empirically (Section 6 of main text), dropping any single term degrades
performance in a characteristic way: without $\mathcal{L}_\FM$ the
support collapses; without $\mathcal{L}_\force$ the local geometry is
noisy at the atomic scale; without $\mathcal{L}_\energy$ the global
Boltzmann shape is not enforced and sample energies have higher
variance than truth.
\end{remark}

\section{Convergence rate}

\begin{theorem}[BGFM convergence in TV]\label{thm:bgfm-conv}
Under Assumptions~\ref{assum:bgfm}, (A1)--(A3) of Appendix A v3, and
assuming the trained model achieves
$\mathcal{L}_\FM \le \varepsilon_\FM$, $\mathcal{L}_\force \le
\varepsilon_\force$, $\mathcal{L}_\energy \le \varepsilon_\energy$,
the retracted-Euler scheme with step $\Delta t$ yields a sample
distribution $p_{\theta^\Delta t}$ with
\[
\TV(p_{\theta^{\Delta t}}, p_\text{Boltz}^\text{data-support})
  \le C_1 \Delta t + C_2 \sqrt{\varepsilon_\FM}
    + C_3 \sqrt{\varepsilon_\force} + C_4 \sqrt{\varepsilon_\energy},
\]
where $p_\text{Boltz}^\text{data-support}$ is the Boltzmann distribution
restricted (normalized) to $\mathrm{supp}(p_\text{data})$, and the
constants depend on Lipschitz constants and the reach $\tau_0$ of
$\M_\ster$.
\end{theorem}

\begin{proof}[Proof sketch]
The TV bound decomposes into four sources:
\begin{enumerate}
\item Euler discretization: $O(\Delta t)$ as in Theorem~7 of v3.
\item FM support error: $O(\sqrt{\varepsilon_\FM})$ by Pinsker
      (controls $\TV$ by $\sqrt{\KL}$, which in turn is bounded by FM
      loss in the straight-line interpolant regime).
\item Force error: a gradient mismatch perturbs the density shape by
      integration: $O(\sqrt{\varepsilon_\force})$.
\item Energy error: direct contribution to the shape: $O(\sqrt{\varepsilon_\energy})$.
\end{enumerate}
Triangle inequality over contributions gives the total bound.
\end{proof}

\begin{corollary}[Sampling rate vs reflected SDE]
Under regularity, BGFM retracted-Euler achieves $\TV$ discretization
error $O(\Delta t)$, whereas reflected-SDE schemes (Lou-Ermon,
Fishman) achieve $O(\sqrt{\Delta t})$. BGFM is thus one order faster
in the discretization regime.
\end{corollary}

\section{Relationship to prior work}

\begin{remark}[Relation to Boltzmann Generators]
Noe et al.~(\textit{Science}, 2019) introduced the idea that a
generative model can be trained to sample from a known Boltzmann
distribution by combining (i) a normalizing-flow density estimator
with (ii) a reweighting (KL) loss. Follow-ups in ML venues
(K\"ohler 2020 ICML, Midgley 2023 ICLR, Falkner 2022 ICML) refined
this for molecular systems.

BGFM differs in three ways:
\begin{enumerate}
\item \textbf{Flow matching, not normalizing flow.} Flow matching
      replaces the bijective-transform density estimator with an ODE,
      yielding $O(\Delta t)$ convergence (vs $O(\sqrt{\Delta t})$ for
      reflected SDE / NF-based approaches) and simpler training.
\item \textbf{Universal neural potential.} Prior Boltzmann generators
      used hand-designed or system-specific potentials (Lennard-Jones,
      harmonic forces). BGFM uses OMol25, a universal pretrained
      neural potential, making the framework applicable to any
      molecular chemistry in a single training.
\item \textbf{Force + energy combined.} Prior works typically combine
      NLL with one reweighting term. BGFM combines flow matching, score
      matching (force), and energy matching simultaneously. Under
      capacity, the three are redundant (Remark after
      Theorem~\ref{thm:boltzmann}), but practically they give
      complementary training signals at distinct scales.
\end{enumerate}
\end{remark}

\begin{remark}[Relation to CDD (Cardei NeurIPS 2025)]
CDD enforces hard combinatorial constraints (e.g., valence) via
differentiable projection onto a discrete constraint set. BGFM is
bond-free and does not use discrete combinatorial projection; chemistry
is instead soft-enforced via the neural potential $E$. The two methods
target orthogonal axes of the constrained-generation problem; they
are complementary rather than competing.
\end{remark}

\bigskip
\hrule
\bigskip

\emph{End of Appendix A v4. See notes/bgfm\_method.md for
implementation details. Formal theorem statements and proof structure
reviewed against Lipman et al.~2023 (FM), Hyv\"arinen 2005 (score
matching), and No\'e et al.~2019 (Boltzmann generator).}

\end{document}
